\section{機械学習の基礎}

機械学習(machine learning)は，観測データから規則性を抽出し，未知の入力に対する予測や意思決定を可能にする方法論である．しかし，その中心課題は単なる経験的な当てはめではない．観測標本を生み出した確率分布を，どのようなモデル族によって近似し，その近似をどのような最適化問題として解くかという二つの問いが，多くの手法の背後にある．本章では，この視点を軸として，損失関数，勾配法，誤差逆伝播法，ニューラルネットワーク，さらに画像・系列・言語を扱う代表的構造へと議論を進める．

本章の目標は三つある．第一に，機械学習を「確率分布の学習と最適化の問題」として統一的に捉えることである．第二に，教師あり学習(supervised learning)を中心として，経験リスク最小化(empirical risk minimization; ERM)と汎化(generalization)の基礎概念を定式化することである．第三に，現代的な深層学習(deep learning)の主要要素である確率的勾配降下法，誤差逆伝播法，畳み込みニューラルネットワーク(convolutional neural network; CNN)，再帰型ニューラルネットワーク(recurrent neural network; RNN)，Transformer が，それぞれどのような問題意識と計算上の制約から導かれたかを理解することである．

\subsection{機械学習の歴史}
機械学習の歴史を学ぶ必要があるのは，現在広く用いられている手法が，単に性能競争の結果として現れたのではなく，表現力，計算資源，データ量，最適化の難しさという四つの制約に対する応答として形成されてきたからである．ある時代の主流手法は，その時点で利用可能な計算機とデータのもとで，どの程度複雑な関数を学習でき，どの程度安定に最適化できるかという条件によって規定されていた．この見取り図を先に持つことは，後に学ぶ CNN，RNN，Transformer を個別の技巧としてではなく，問題構造に応じた帰結として理解するために重要である．本章の議論は \citet{Bishop2006}，\citet{Murphy2012}，\citet{Goodfellow2016} などの標準的教科書を踏まえつつ，本書の関心に沿って構成する．

初期のパターン認識(pattern recognition)では，主たる課題は観測ベクトルからクラスを識別する判別規則を設計することであった．ここでは，特徴抽出は人手で設計されることが多く，識別器は線形判別，最近傍法，判別分析など比較的単純な形をとった．この段階では，モデルの表現力は限定的であったが，そのぶん計算量は制御しやすく，少量のデータでも扱いやすかった．他方で，入力表現そのものを学習する能力が弱く，複雑な知覚課題には限界があった．

これに対して統計的学習(statistical learning)\cite{Vapnik1998,Hastie2009}は，有限標本からの推定と汎化誤差の制御を理論的に扱う枠組みを整備した．経験リスク最小化，正則化，構造的リスク最小化，カーネル法などは，モデルの容量と汎化性能の折り合いをどのように付けるかという問いに答えるものである．この流れでは，表現力の拡張と同時に，過学習を避けるための理論的統制が重視された．すなわち，機械学習は単に当てはまりのよい関数を探す問題ではなく，未知データに対しても安定に振る舞う関数を推定する問題であるという認識が明確になった．

ニューラルネットワーク(neural network)の発想自体は古いが，1980年代に誤差逆伝播法が広まったことで，多層モデルの学習が現実的な手続きとして定着した．ここで重要なのは，ニューラルネットワークが高い表現力を持ちながら，同時に最適化の難しさを強く抱え込む点である．浅いモデルでも非線形関数近似は可能であったが，計算資源が限られ，データ量も小さく，初期化や活性化関数の設計も未成熟であったため，深いモデルは長らく安定に学習しにくかった．したがって，表現力だけでは実用的成功は保証されず，それを支える最適化技術と計算基盤が不可欠であった．

2000年代後半以降の深層学習の成功は，この四つの制約が同時に緩和されたことによって説明できる．第一に，GPU などの計算資源の発展により，大規模な行列演算を反復する訓練が実用化した．第二に，Web やセンサの普及によって大規模データが利用可能になった．第三に，ReLU 系活性化関数，適切な初期化，正規化，SGD の改良などにより，深いネットワークの最適化が容易になった．第四に，深い階層構造が画像・音声・言語における多段の表現学習に適していることが経験的に確認された．この結果，特徴量設計の重心が人手から学習へと移った．

現代の基盤モデル(foundation model)や大規模モデル(large-scale model)は，この延長線上に位置する．大規模事前学習によって汎用的な表現を獲得し，その後に少量の追加データや指示で多様な下流課題へ適応するという発想は，モデルの表現力とデータ量を極限まで押し広げた結果である．ただし，その基盤にあるのは依然として，条件付き分布や系列分布の学習を，巨大な最適化問題として解くという構図である．

この歴史的流れを各構造に結びつけると，CNN は画像に固有な局所性と平行移動対称性を取り込むために現れ，RNN は系列データの時間依存を表現するために導入され，Transformer は RNN の逐次性と長距離依存性の弱さを克服するために発展した．したがって，以後の節で扱う個々のモデルは，すべて「どのような構造仮定を置けば，限られた計算のもとで確率分布をうまく学習できるか」という共通問題に対する別々の答えである．次節では，この問題を教師あり学習の基本設定として明示的に定式化する．

\subsection{機械学習の基本設定}
機械学習の議論を精密に進めるためには，まず何を観測し，何を予測し，何を最適化するのかを明確にしなければならない．直感的には，入力から出力への対応を学ぶことが目的であるが，その背後では未知のデータ生成分布のもとで予測誤差を小さくするという問題が立てられている．とりわけ教師あり学習では，入力と正解出力が対になった標本が与えられるため，確率的定式化と最適化問題との対応がもっとも見えやすい．

\paragraph*{定義}
教師あり学習では，訓練データセットを
\begin{equation}
\mathcal D = \{(x_i,y_i)\}_{i=1}^n
\end{equation}
と書く．ここで，$x_i \in \mathcal X$ は第 $i$ 標本の入力，$y_i \in \mathcal Y$ は対応する出力，$n$ は標本数，$\mathcal X$ と $\mathcal Y$ はそれぞれ入力空間と出力空間である．通常は，各対 $(x_i,y_i)$ が未知の同時分布 $p_{\mathrm{data}}(x,y)$ から独立同分布に従って生成されたと仮定する．この仮定は理想化ではあるが，訓練データと将来のデータが同種の母集団から来るという汎化の前提を数学的に表現している．

学習器はパラメータ(parameter)$\theta$ を持つモデル $f_\theta$ として表される．$f_\theta(x)$ は入力 $x$ に対する予測値であり，その形は問題設定に応じて異なる．回帰(regression)では $\mathcal Y$ は実数または実ベクトルであり，$f_\theta(x)$ は平均値や連続量の予測を返す．分類(classification)では $\mathcal Y$ は有限個のクラス集合であり，$f_\theta(x)$ はクラス確率やロジット(logit)を返すことが多い．いずれにせよ，$\theta$ はモデルの自由度を定める未知量であり，学習とは $\theta$ をデータから推定することである．

実際の学習では，データは訓練集合，検証集合，テスト集合に分割される．訓練集合はパラメータ更新に用いられ，検証集合はハイパーパラメータの選択や早期終了の判断に用いられ，テスト集合は最終的な汎化性能の評価に用いられる．この分割が必要なのは，同じデータでモデルの選択と性能評価を同時に行うと，評価が楽観的に歪むからである．とくに深層学習では，アーキテクチャ，学習率，正則化強度など多くの設計選択があるため，検証集合の役割は本質的である．

教師あり学習の中心概念は経験リスク最小化である．損失関数(loss function)を $\ell(y,f_\theta(x))$ と書くと，経験リスクは
\begin{equation}
\hat R_n(\theta)= \frac{1}{n}\sum_{i=1}^n \ell(y_i,f_\theta(x_i))
\end{equation}
で定義される．ここで $\ell$ は予測 $f_\theta(x_i)$ が正解 $y_i$ からどれだけ外れているかを測る関数であり，$\hat R_n(\theta)$ は訓練データ上の平均誤差である．他方，真に最小化したい量は未知分布に関する期待損失
\begin{equation}
R(\theta)= \mathrm{E}_{(x,y)\sim p_{\mathrm{data}}}\left[\ell(y,f_\theta(x))\right]
\end{equation}
であり，これを期待リスクまたは汎化誤差と呼ぶ．経験リスク最小化は，観測できない $R(\theta)$ を，その標本平均である $\hat R_n(\theta)$ に置き換えて最適化する考え方である．

このとき重要なのは，モデルが単に点予測関数としてではなく，条件付き分布 $p(y|x)$ の近似として理解できることである．すなわち，$f_\theta(x)$ はしばしば直接の予測値ではなく，条件付き分布 $p_\theta(y|x)$ を指定するためのパラメータを返す．たとえば回帰では，$f_\theta(x)$ を条件付き平均 $\mu_\theta(x)$ とみなして
$p_\theta(y|x)$ をガウス分布で表すことが多い．分類では，$f_\theta(x)$ がクラスごとのロジットを返し，それを softmax に通してクラス確率 $p_\theta(y=k|x)$ を得る．この対応により，損失関数は確率モデルの負の対数尤度として自然に導かれる．

さらに，損失関数と条件付き分布の関係を理解すると，最適予測器の意味も明確になる．たとえば二乗損失のもとでは，各入力 $x$ に対する最適予測は条件付き期待値 $\mathrm{E}[Y|X=x]$ である．分類で $0$-$1$ 損失を考えれば，最適予測は事後確率 $p(y|x)$ が最大のクラスを返す規則になる．したがって，機械学習の目標は単なるラベルの記憶ではなく，条件付き分布の構造を十分によく近似し，適切な決定規則へ変換することである．

以上の定式化により，問題は「どの損失関数を選び，どのような確率モデルを仮定するか」に帰着する．回帰と分類が異なるのは出力空間の性質だけではなく，仮定する確率モデルと対応する損失が異なるためである．次節では，この対応を最尤推定と最大事後確率推定の観点から詳しく導き，損失最小化がどのような最適化問題になるかを整理する．

\subsection{損失関数と最適化}
損失関数を導入する理由は，学習の目的を数理的に定義し，パラメータ推定を最適化問題へ変換するためである．経験的には二乗誤差やクロスエントロピーがよく用いられるが，これらは便宜的に選ばれているわけではない．適切な確率モデルを仮定すると，それらは最尤推定(maximum likelihood estimation; MLE)や最大事後確率推定(maximum a posteriori; MAP)の帰結として現れる．したがって，損失の選択は「どの誤差を罰したいか」という感覚的判断だけではなく，「データ生成過程をどのようにモデル化するか」という統計的判断でもある．

\paragraph*{定義}
最尤推定では，パラメータ $\theta$ を，訓練データの条件付き尤度を最大化する量として定める．すなわち
\begin{equation}
\hat\theta_{\mathrm{MLE}}
= \arg\max_\theta \prod_{i=1}^n p_\theta(y_i|x_i)
= \arg\max_\theta \sum_{i=1}^n \log p_\theta(y_i|x_i)
\end{equation}
である．対数は単調増加であるため，尤度最大化は負の対数尤度の最小化と等価である．以後，負の対数尤度を損失関数として用いると，確率推定と最適化が一つの枠組みに統合される．

回帰において，出力 $y_i$ が平均 $\mu_\theta(x_i)$，分散 $\sigma^2$ のガウス分布に従うと仮定しよう．ここで $\mu_\theta(x_i)$ はモデルが入力 $x_i$ から出力する条件付き平均，$\sigma^2$ は観測雑音の分散である．このとき
\begin{equation}
p_\theta(y_i|x_i)
= \frac{1}{\sqrt{2\pi\sigma^2}}
\exp\left(
-\frac{(y_i-\mu_\theta(x_i))^2}{2\sigma^2}
\right)
\end{equation}
であるから，負の対数尤度の総和は定数項を除いて
\begin{equation}
\sum_{i=1}^n -\log p_\theta(y_i|x_i)
= \frac{1}{2\sigma^2}\sum_{i=1}^n(y_i-\mu_\theta(x_i))^2 + \mathrm{const.}
\end{equation}
となる．したがって，平均二乗誤差(mean squared error; MSE)
\begin{equation}
\mathcal L_{\mathrm{MSE}}(\theta)
= \frac{1}{n}\sum_{i=1}^n(y_i-\mu_\theta(x_i))^2
\end{equation}
は，ガウス雑音仮定のもとでの最尤推定に対応している．ここで $n$ は標本数であり，各誤差を平均することでデータ数に依らない尺度を得ている．

分類では，出力が $K$ 個のクラスのいずれかをとるとする．モデルが入力 $x_i$ に対してクラスごとのロジット $z_k(x_i)$ を返すとき，softmax により
\begin{equation}
\pi_{\theta,k}(x_i)
= \frac{\exp(z_k(x_i))}{\sum_{j=1}^K \exp(z_j(x_i))}
\end{equation}
を定義すると，$\pi_{\theta,k}(x_i)$ はクラス $k$ の条件付き確率 $p_\theta(y_i=k|x_i)$ と解釈できる．正解ラベルを one-hot ベクトル $y_i=(y_{i1},\dots,y_{iK})$ で表し，$y_{ik}\in\{0,1\}$ を「標本 $i$ がクラス $k$ に属するか」を表す指示変数とすると，負の対数尤度は
\begin{equation}
\mathcal L_{\mathrm{CE}}(\theta)
= -\sum_{i=1}^n \sum_{k=1}^K y_{ik}\log \pi_{\theta,k}(x_i)
\end{equation}
となる．これはクロスエントロピー損失(cross-entropy loss)であり，真のラベル分布とモデル分布のずれを測る量として理解できる．

MAP 推定では，さらにパラメータの事前分布 $p(\theta)$ を導入し，
\begin{equation}
\hat\theta_{\mathrm{MAP}}
= \arg\max_\theta
\left\{
\sum_{i=1}^n \log p_\theta(y_i|x_i)+ \log p(\theta)
\right\}
\end{equation}
を解く．これは負の対数をとれば，データ適合項に加えて $-\log p(\theta)$ という正則化項を課すことに等しい．たとえば $p(\theta)\propto \exp(-\lambda \|\theta\|_2^2/2)$ というガウス事前分布を置けば，$\lambda \|\theta\|_2^2/2$ という $L_2$ 正則化が現れる．同様に，ラプラス事前分布を置けば $L_1$ 正則化が現れる．したがって，正則化は経験的な細工ではなく，パラメータに対する事前信念を導入した統計推定として理解できる．

もっとも，損失関数が定まれば問題が簡単になるわけではない．線形回帰やロジスティック回帰の一部では凸最適化\cite{Boyd2004}が得られるが，深いニューラルネットワークでは一般に目的関数は非凸(nonconvex)である．すなわち，勾配がゼロとなる停留点の周囲には，局所最小点，鞍点(saddle point)，あるいは曲率が小さい平坦領域が存在しうる．ヘッセ行列を $H(\theta)$ と書けば，停留点で $H(\theta)$ の固有値がすべて正なら局所最小，正負が混在すれば鞍点である．深層学習においては，悪い局所解そのものよりも，高次元空間に多数現れる鞍点や，曲率の偏りによって進行が遅くなる平坦領域のほうが，しばしば実践上の難所となる．

さらに，最小化された訓練損失が小さいことと，汎化性能が高いことは同義ではない．平坦な最小点はしばしば汎化に有利だと経験的に言われるが，その因果関係は単純ではなく，パラメータ化の仕方や最適化アルゴリズムにも依存する．ここで重要なのは，機械学習における最適化は，解析的に唯一解を求める問題ではなく，巨大な非凸目的関数の中から，計算可能で汎化も良好な解を探索する過程だという点である．

この探索を現実的な時間で実行するために，深層学習では勾配法が中心的役割を果たす．とくに大規模データでは，全データに対する勾配を毎回計算することは高価すぎる．次節では，この問題に対する標準的解法である確率的勾配降下法を導入し，その統計的性質と実装上の意味を整理する．

\subsection{確率的勾配降下法(SGD)}
損失関数が定まった後に必要なのは，その最小値へ向かってパラメータを更新する具体的な計算法である．とくに深層学習では，パラメータ数もデータ数も大きいため，厳密解を求めるよりも反復的な近似最適化が現実的である．このとき，全訓練データを毎回用いる方法は統計的には素直であっても，計算量とメモリ使用量の点でしばしば実用的でない．確率的勾配降下法(stochastic gradient descent; SGD)は，この制約のもとで勾配情報を効率よく利用するための基本アルゴリズムであり，その源流は\citet{RobbinsMonro1951}の確率近似法にさかのぼる．

訓練目的関数を
\begin{equation}
L(\theta)=\frac{1}{n}\sum_{i=1}^n \ell_i(\theta),
\qquad
\ell_i(\theta)=\ell(y_i,f_\theta(x_i))
\end{equation}
と書く．ここで $L(\theta)$ は平均損失，$\ell_i(\theta)$ は標本 $i$ に対する個別損失，$n$ は訓練標本数である．バッチ勾配降下法(batch gradient descent)では，各反復 $t$ における更新は
\begin{equation}
\theta_{t+1}
= \theta_t - \eta_t \nabla L(\theta_t)
= \theta_t - \eta_t \frac{1}{n}\sum_{i=1}^n \nabla \ell_i(\theta_t)
\end{equation}
で与えられる．ここで $\theta_t$ は時刻 $t$ のパラメータ，$\eta_t>0$ は学習率(learning rate)である．この更新は目的関数の真の勾配を用いるため安定であるが，大規模データでは毎回すべての標本を読む必要がある．

これに対して SGD では，訓練集合からサイズ $b$ のミニバッチ $B_t$ を無作為に選び，
\begin{equation}
g_t
= \frac{1}{b}\sum_{i\in B_t}\nabla \ell_i(\theta_t),
\qquad
\theta_{t+1}=\theta_t-\eta_t g_t
\end{equation}
と更新する．ここで $b$ はバッチサイズ，$g_t$ はミニバッチ勾配である．$B_t$ を一様にサンプリングすれば，条件付き期待値について
\begin{equation}
\mathrm{E}[g_t\mid \theta_t] = \nabla L(\theta_t)
\end{equation}
が成り立つ．すなわち，ミニバッチ勾配は真の勾配の不偏推定量である．この事実が，SGD を単なる近似計算ではなく，統計的に意味のある勾配推定法として位置づけている．

ミニバッチ化の効果を理解するには，勾配の分散に注目するとよい．個別勾配 $\nabla \ell_i(\theta_t)$ がある共分散行列 $\Sigma(\theta_t)$ を持つと近似すれば，復元抽出に基づく単純化のもとで
\begin{equation}
\mathrm{Var}(g_t\mid \theta_t)\approx \frac{1}{b}\Sigma(\theta_t)
\end{equation}
である．したがって，バッチサイズ $b$ を大きくすると勾配推定の分散はおおむね $1/b$ に比例して減少する．これは，大きなバッチでは更新方向が安定し，小さなバッチでは雑音が強くなることを意味する．雑音は収束を遅らせうる一方で，鞍点や鋭い谷から抜け出す助けになる場合もあるため，バッチサイズの選択は統計的効率と探索性の折り合いである．

実装上，ミニバッチ化が重要なのは，第一にメモリ制約のためであり，第二に計算資源の並列利用のためである．GPU や TPU は多数の標本に対する同型演算をまとめて実行することで高いスループットを発揮するが，全データを一度に載せることは通常できない．ミニバッチは，行列演算の効率を確保しつつ，メモリ内で処理可能な単位にデータを分割する役割を持つ．また，オンラインに到着するデータや巨大データセットに対しても，局所的な勾配情報だけで更新を続けられる点で，大規模学習の前提をなしている．

\paragraph*{アルゴリズム}
SGD の基本形は単純である．各エポックでデータをシャッフルし，ミニバッチごとに順伝播，損失計算，勾配計算，パラメータ更新を行う．ここで「エポック」とは訓練集合全体を一巡する単位を指す．実際には，学習率の減衰，勾配クリッピング，重み減衰などを組み合わせることが多いが，中心にあるのは不偏な勾配推定を反復的に用いて目的関数を下げるという構図である．

SGD の派生として Momentum は過去の更新方向を指数移動平均し，振動を抑えながら進行方向を安定化する．Adam は一次モーメントと二次モーメントの推定に基づき，各パラメータに適応的な学習率を割り当てる．これらは実務上きわめて有用であるが，深層学習の根幹がミニバッチ勾配に基づく SGD にあることは変わらない．勾配法が機能する前提は，各反復でパラメータに関する勾配を効率よく計算できることであり，そのための決定的技術が誤差逆伝播法である．次節では，連鎖律に基づく勾配計算の仕組みを導く．

\subsection{誤差逆伝播法}
深いネットワークを学習するには，多数のパラメータそれぞれについて損失の勾配を求める必要がある．もし各パラメータを独立に微分し，そのたびに順方向計算をやり直すなら，計算量は現実的でなくなる．誤差逆伝播法(backpropagation)\cite{Rumelhart1986}が必要なのは，複合関数として構成されたモデルの勾配を，計算グラフと連鎖律を用いて一度の順伝播と一度の逆伝播で効率よく求めるためである．この方法により，パラメータ数が多くても，計算量はグラフの辺数にほぼ線形に比例する．

計算グラフ(computational graph)では，各演算をノード，依存関係を辺として表す．$L$ 層の全結合ネットワークを例にとると，入力を $h^{(0)}=x$ として，各層 $l=1,\dots,L$ について
\begin{equation}
a^{(l)} = W^{(l)} h^{(l-1)} + b^{(l)},
\qquad
h^{(l)} = \phi^{(l)}(a^{(l)})
\end{equation}
と書ける．ここで $W^{(l)}$ は第 $l$ 層の重み行列，$b^{(l)}$ はバイアスベクトル，$a^{(l)}$ は活性化関数適用前の線形結合，$h^{(l)}$ は層出力，$\phi^{(l)}$ は活性化関数である．教師信号を $y$，最終損失を $\mathcal L(h^{(L)},y)$ と書けば，求めたいのは各 $W^{(l)}$ と $b^{(l)}$ に関する偏微分である．

\paragraph*{定義}
逆伝播では，各層の前活性化 $a^{(l)}$ に関する損失の勾配
\begin{equation}
\delta^{(l)} = \frac{\partial \mathcal L}{\partial a^{(l)}}
\end{equation}
を誤差信号と呼ぶ．出力層では
\begin{equation}
\delta^{(L)}
= \frac{\partial \mathcal L}{\partial h^{(L)}} \odot \phi^{(L)\prime}(a^{(L)})
\end{equation}
であり，$\odot$ は要素ごとの積，$\phi^{(L)\prime}$ は活性化関数の導関数である．中間層では連鎖律により
\begin{equation}
\delta^{(l)}
= \left((W^{(l+1)})^\top \delta^{(l+1)}\right)\odot \phi^{(l)\prime}(a^{(l)}),
\qquad l=L-1,\dots,1
\end{equation}
という再帰式が成り立つ．すなわち，上位層から伝わってきた誤差を重み行列の転置で引き戻し，その層の局所微分で修正することで，下位層の勾配情報が得られる．

この誤差信号が得られれば，各パラメータに関する勾配は直ちに計算できる．単一標本に対しては
\begin{equation}
\frac{\partial \mathcal L}{\partial W^{(l)}}
= \delta^{(l)}(h^{(l-1)})^\top,
\qquad
\frac{\partial \mathcal L}{\partial b^{(l)}}
= \delta^{(l)}
\end{equation}
である．ミニバッチ学習では，各標本に対するこれらの勾配を平均したものがバッチ勾配となる．ここで $h^{(l-1)}$ が必要なのは，重み行列 $W^{(l)}$ が前層出力と線形に結びついているからであり，誤差信号 $\delta^{(l)}$ と入力 $h^{(l-1)}$ の外積として勾配が現れる．

順伝播時の中間値を保存する必要があるのも，この式から理解できる．逆伝播では $\phi^{(l)\prime}(a^{(l)})$ や $h^{(l-1)}$ が必要になるため，前活性化 $a^{(l)}$ と層出力 $h^{(l)}$ を保持しておかないと，逆向き計算のたびに順方向を再計算しなければならない．したがって，誤差逆伝播法は計算量を節約する代わりに，中間活性を記憶するメモリを要求する．この時間と空間の交換は，深いモデルや長い系列を扱う際にとくに重要である．

逆伝播が効率的である理由は，逆向きモード自動微分(reverse-mode automatic differentiation)と同じ原理を用いているからである．出力がスカラー損失一つであるとき，各ノードから損失への感度を一度だけ伝播すれば，すべてのパラメータの勾配が得られる．これに対して前向きモードで各パラメータごとに微分を流すと，パラメータ数に比例する回数の計算が必要になる．深層学習では通常，出力次元よりもパラメータ数のほうがはるかに大きいため，逆伝播の利点は決定的である．

もっとも，深いネットワークでは勾配そのものが安定に伝わらないことがある．層ごとのヤコビアンを $J^{(l)}$ と書けば，下位層への勾配は概ねそれらの積
\begin{equation}
\frac{\partial \mathcal L}{\partial h^{(l)}}
\propto
J^{(l+1)}J^{(l+2)}\cdots J^{(L)}
\end{equation}
に支配される．この積のノルムが繰り返し $1$ より小さければ勾配消失(vanishing gradient)，大きければ勾配爆発(exploding gradient)が起こる．深層化と長系列化に伴ってこの問題は顕著になり，活性化関数，初期化，正規化，残差接続，ゲーティングなどの工夫が必要になる．

このように，誤差逆伝播法は深層学習を可能にした計算技術であるが，その効果はネットワークの構造と強く結びついている．次節では，そもそもニューラルネットワークがどのような関数族を表し，なぜ非線形性が不可欠なのかを整理する．

\subsection{ニューラルネットワークの基本構造}
ニューラルネットワークを学ぶ必要があるのは，深層学習におけるほとんどすべてのモデルが，その変形または拡張として理解できるからである．CNN，RNN，Transformer はいずれも特殊な構造を持つが，その基本単位は線形変換と非線形活性化の反復にある．したがって，まずパーセプトロン，多層パーセプトロン，活性化関数，出力層の意味を押さえることが，後続の複雑なアーキテクチャを理解する前提となる．

パーセプトロン(perceptron)は，もっとも基本的な線形判別器である．入力 $x\in\mathcal X$ をベクトルとし，重みベクトルを $w$，バイアスを $b$ とすると，その決定関数は
\begin{equation}
\hat y = \mathrm{sign}(w^\top x + b)
\end{equation}
と書ける．ここで $w^\top x$ は入力と重みの内積，$\mathrm{sign}$ は符号関数であり，超平面による二値分類を表す．パーセプトロンは線形分離可能な問題には有効であるが，入力空間における複雑な非線形境界を直接表すことはできない．この限界を超えるために，中間表現を導入した多層パーセプトロン(multilayer perceptron; MLP)が用いられる．

MLP では，各層が
\begin{equation}
h^{(l)} = \phi^{(l)}(W^{(l)}h^{(l-1)} + b^{(l)})
\end{equation}
という形で定義される．ここで $h^{(0)}=x$ は入力，$h^{(l)}$ は第 $l$ 層の隠れ表現，$W^{(l)}$ と $b^{(l)}$ は学習されるパラメータ，$\phi^{(l)}$ は活性化関数である．最終層の出力は課題に応じて変わる．回帰では線形出力をそのまま使うことが多く，二値分類では sigmoid 関数で確率に写し，多クラス分類では softmax を用いてクラス確率へ変換する．

非線形性が必要であることは，数式から直ちに分かる．もしすべての活性化関数が恒等写像，すなわち $\phi^{(l)}(u)=u$ であるなら，二層以上のネットワークでも
\begin{equation}
h^{(L)} = W_{\mathrm{eff}}x + b_{\mathrm{eff}}
\end{equation}
という単一のアフィン変換にまとめられる．ここで $W_{\mathrm{eff}}$ は各層の重み行列の積，$b_{\mathrm{eff}}$ はバイアス項をまとめたベクトルである．したがって，線形変換を何層重ねても，表現力は本質的に一層の線形モデルを超えない．非線形活性化は，関数クラスを真に拡張し，高次の特徴や複雑な決定境界を表現するために不可欠である．

活性化関数にはそれぞれ役割と副作用がある．sigmoid 関数
\begin{equation}
\sigma(u)=\frac{1}{1+\exp(-u)}
\end{equation}
は出力を $(0,1)$ に収めるため確率解釈と相性がよいが，大きな正負の入力で勾配が小さくなりやすい．双曲線正接関数 $\tanh(u)$ は出力が $(-1,1)$ に中心化されるため学習しやすい場面があるが，やはり飽和による勾配消失を起こしうる．ReLU は
\begin{equation}
\mathrm{ReLU}(u)=\max(0,u)
\end{equation}
で定義され，正の領域で勾配が保たれるため深層学習を大きく前進させた．近年の Transformer では，標準正規分布の累積分布関数を $\Phi$ として $\mathrm{GELU}(u)=u\Phi(u)$ と定義される GELU がよく用いられる．これは入力を滑らかにゲートする働きを持ち，深いモデルの最適化と相性がよい．

普遍近似定理(universal approximation theorem)は，適切な非線形活性化を持つ十分広い一隠れ層ネットワークが，コンパクト集合上の連続関数を任意精度で近似できることを述べる．この結果はニューラルネットワークの表現力を保証するが，それだけで深いネットワークの優位性が自動的に示されるわけではない．実際には，同じ関数を浅いネットワークで表すには非常に大きな幅が必要になることがあり，深さは合成構造を持つ関数を効率よく表す手段として重要である．すなわち，深いモデルの意義は「表現できるか否か」だけでなく，「どれだけ少ないパラメータで，どれだけ学習しやすく表現できるか」にある．

以上より，ニューラルネットワークの基本単位は，線形写像，非線形活性化，目的に応じた出力層から成ることが分かる．しかし高い表現力は同時に過学習の危険を伴う．次節では，訓練誤差の最小化と未知データへの汎化との緊張関係を整理し，正則化の役割を考える．

\subsection{正則化と汎化}
機械学習で本当に必要なのは，訓練データによく合うモデルではなく，未知データに対しても安定に機能するモデルである．正則化(regularization)を学ぶ理由は，モデルの表現力が高くなるほど，訓練誤差を小さくするだけでは不十分になるからである．とくに深層学習では，過剰な自由度が過学習を招く一方で，十分な表現力がなければ複雑な分布を捉えられない．正則化はこの緊張関係を制御し，汎化性能を改善するための原理である．

過学習(overfitting)とは，モデルが訓練集合に特有の偶然的変動や雑音まで取り込んでしまい，訓練誤差は小さいにもかかわらずテスト誤差が大きくなる現象をいう．逆に，モデルが単純すぎて訓練誤差すら十分に下がらない状態をアンダーフィット(underfitting)という．前者ではモデル容量がデータ量に比して大きすぎ，後者では表現力が不足している．したがって，良い学習とは，単に大きなモデルを選ぶことでも，小さなモデルを選ぶことでもなく，課題に見合った容量制御を行うことである．

この関係を理解する一つの手がかりが，二乗損失におけるバイアス・バリアンス分解(bias-variance decomposition)である．真の回帰関数を $f^\star(x)=\mathrm{E}[Y|X=x]$，学習データ $\mathcal D$ から得られる推定器を $\hat f_{\mathcal D}(x)$，観測雑音の分散を $\sigma_\varepsilon^2=\mathrm{Var}(Y|X=x)$ とすると，固定した入力 $x$ における予測誤差の期待値は
\begin{equation}
\mathrm{E}_{\mathcal D,Y}\left[(Y-\hat f_{\mathcal D}(x))^2\right]
= \sigma_\varepsilon^2
+ \left(\mathrm{E}_{\mathcal D}[\hat f_{\mathcal D}(x)]-f^\star(x)\right)^2
+ \mathrm{E}_{\mathcal D}\left[
\left(
\hat f_{\mathcal D}(x)-\mathrm{E}_{\mathcal D}[\hat f_{\mathcal D}(x)]
\right)^2
\right]
\end{equation}
と書ける．右辺の第一項は不可避な雑音，第二項は系統的なずれであるバイアス，第三項は学習データの揺らぎに対する感度であるバリアンスを表す．一般に，単純なモデルは高バイアス・低バリアンス，複雑なモデルは低バイアス・高バリアンスの傾向を持つ．

もっとも基本的な正則化は，目的関数にパラメータノルムの罰則を加える方法である．$L_2$ 正則化は $\lambda\|\theta\|_2^2/2$ を，$L_1$ 正則化は $\lambda\|\theta\|_1$ を損失に加える．ここで $\lambda>0$ は正則化強度であり，$\|\theta\|_2$ はユークリッドノルム，$\|\theta\|_1$ は各成分の絶対値和である．前節で見たように，これらはそれぞれガウス事前分布とラプラス事前分布を課した MAP 推定として解釈できる．$L_2$ はパラメータを滑らかに縮小し，$L_1$ は疎な解を促しやすい．いずれも，無制約に複雑な関数を選ぶことを抑える容量制御とみなせる．

深層学習では，明示的な罰則以外にも多くの正則化手段が使われる．Dropout\cite{Srivastava2014} では，訓練時に各隠れユニットを確率的に無効化し，ネットワークが特定の経路に過度に依存することを防ぐ．記号的には，隠れ表現 $h$ に対してベルヌーイマスク $m$ を掛けた $\tilde h = m\odot h$ を用いる操作とみなせる．Early Stopping は，検証誤差が悪化し始める前に訓練を止める方法であり，反復最適化それ自体を暗黙の正則化として利用する．Batch Normalization\cite{Ioffe2015} は，ミニバッチ平均 $\mu_B$ と分散 $\sigma_B^2$ を用いて
\begin{equation}
\hat a = \frac{a-\mu_B}{\sqrt{\sigma_B^2+\epsilon}},
\qquad
y = \gamma \hat a + \beta
\end{equation}
と正規化する．ここで $a$ は正規化前の活性，$\epsilon$ は数値安定化のための小定数，$\gamma,\beta$ は学習可能なスケールとシフトである．これは主として最適化を安定化する技法であるが，ミニバッチ統計量の揺らぎを通じて一定の正則化効果も持つ．

Data Augmentation は，入力にラベル不変な変換を加えることで，モデルに望ましい不変性を学習させる方法である．画像であれば平行移動，反転，切り出し，色変換などが典型であり，系列や言語でもノイズ付与やマスク化などが利用される．これは，単に標本数を人工的に増やす手法ではなく，「この変換では意味が変わらない」という事前知識を学習問題に埋め込む操作と解釈できる．その意味で，アーキテクチャ設計と同様に，分布に関する帰納バイアス(inductive bias)を与える手段である．

近年は，モデルサイズを増やすとテスト誤差が一度悪化した後で再び改善する double descent も観測されている．これは古典的な「モデルが複雑になるほど汎化は単調に悪化する」という単純図式が十分でないことを示している．しかし，この現象もまた，補間可能な過剰パラメータ領域であっても，どの解が最適化によって選ばれるかが重要であることを示すにすぎない．すなわち，汎化は依然として，データ，モデル容量，最適化，正則化の相互作用として理解すべき問題である．

ここまでの議論は，主として全結合ネットワークを念頭に置いていた．しかし現実のデータには，画像の局所性や系列の時間依存のような明確な構造が存在する．こうした構造を活用することで，表現効率と汎化性能を同時に高めることができる．次節では，画像に対する代表的な帰納バイアスを実装した CNN を扱う．

\subsection{CNN}
画像に対して全結合層だけを用いることが非効率である理由は，画像が高次元でありながら，局所的で平行移動に対して類似した構造を持つからである．たとえば，高さ $H$，幅 $W$，チャネル数 $C$ の画像を一つのベクトルに平坦化して全結合層へ入力すると，隠れ層のユニット数に比例して巨大な数のパラメータが必要になる．しかも，この表現では近くにある画素と遠くにある画素の区別が消え，画像の空間構造が活かされない．CNN は，この問題を局所結合と重み共有によって解決する．

畳み込み層では，各出力ユニットが入力全体ではなく局所領域だけを見る．これを受容野(receptive field)と呼ぶ．さらに，入力の異なる位置で同じカーネルを使い回すことで重み共有を行う．二次元入力 $X$ とカーネル $K$ に対する離散畳み込みを，出力位置 $(i,j)$ で
\begin{equation}
(X*K)_{i,j}
= \sum_{u=0}^{r-1}\sum_{v=0}^{s-1} K_{u,v}X_{i+u,j+v}
\end{equation}
と定義する．ここで $r,s$ はカーネルの高さと幅，$u,v$ はカーネル内部の添字，$i,j$ は出力位置を表す．実際の画像ではチャネル方向の和も加わるが，本質は局所領域に対して同一の線形演算を空間全体で適用する点にある．

この構造は二つの利点をもたらす．第一に，パラメータ数が大幅に削減される．全結合層では入力次元に比例して重み数が増えるのに対し，畳み込み層ではカーネルサイズに比例するだけである．第二に，平行移動等変性(translation equivariance)が得られる．平行移動演算を $T_\Delta$，畳み込み演算を $\mathcal C$ と書けば，
\begin{equation}
\mathcal C(T_\Delta X)= T_\Delta(\mathcal C(X))
\end{equation}
が成り立つ．すなわち，入力をずらせば出力も同じだけずれる．この性質により，同じ局所パターンを画像中のどこでも検出できる．

プーリング(pooling)は，局所領域の代表値をとることで空間解像度を下げ，わずかな位置ずれに対する不変性を高める操作である．最大値をとる max pooling と平均をとる average pooling が典型である．畳み込みとプーリングを重ねると，浅い層ではエッジや局所テクスチャ，深い層ではより抽象的な形状や部品が表現されるようになる．また，層を重ねるにつれて受容野が広がるため，局所情報から大域的文脈へと表現が統合されていく．

CNN の歴史的発展も，構造仮定と計算資源の相互作用として理解できる．LeNet\cite{LeCun1998} は文字認識で畳み込みの有効性を示した初期例であり，AlexNet\cite{Krizhevsky2012} は GPU と ReLU を活用して ImageNet で深層 CNN の優位を広く認知させた．VGG は小さなカーネルを多層に積む設計で深さの効果を明確にし，ResNet\cite{He2016} は
\begin{equation}
h^{(l+1)} = h^{(l)} + F(h^{(l)})
\end{equation}
という残差接続によって，より深いネットワークの最適化を可能にした．ここで $F$ は残差写像であり，恒等写像を近道として残すことで勾配伝播を助ける．

もっとも，CNN の帰納バイアスは画像のような格子状データには適していても，系列のように順序依存が本質となる問題にはそのままでは十分でない．時系列，文章，音声のようなデータでは，過去の情報をどのように状態として保持し，将来の予測に反映させるかが中心問題となる．次節では，そのために導入された RNN を扱う．

\subsection{RNN}
系列データを扱うためには，観測が独立ではなく，時間や位置に沿って依存していることを明示的に表現しなければならない．文章，音声，交通流時系列などでは，現在の観測や予測が過去の文脈に依存するため，入力を単に固定長ベクトルとして扱うだけでは不十分である．RNN が必要なのは，系列の結合分布を連鎖律に従って分解し，その文脈を隠れ状態として逐次的に要約できるからである．

長さ $T$ の系列 $x_{1:T}=(x_1,\dots,x_T)$ の結合分布は，確率の連鎖律により
\begin{equation}
p(x_{1:T})
= \prod_{t=1}^T p(x_t\mid x_{1:t-1})
\end{equation}
と書ける．ここで $x_{1:t-1}$ は時刻 $1$ から $t-1$ までの過去系列である．条件付き生成や系列ラベリングでも同様に，現在の出力は過去の入力や過去の出力に依存する．この依存構造を有限次元の状態で近似するため，RNN は各時刻 $t$ に隠れ状態 $h_t$ を持ち，
\begin{equation}
h_t = \phi(W_{xh}x_t + W_{hh}h_{t-1} + b_h)
\end{equation}
と更新する．ここで $x_t$ は時刻 $t$ の入力，$h_t$ はその時点までの文脈を要約した隠れ状態，$W_{xh}$ と $W_{hh}$ は入力と前時刻状態に対する重み行列，$b_h$ はバイアス，$\phi$ は活性化関数である．出力 $o_t$ は通常
\begin{equation}
o_t = W_{hy}h_t + b_y
\end{equation}
とし，必要に応じて softmax などへ通す．

この再帰式の直感は明快である．各時刻で新しい観測 $x_t$ を取り込み，それ以前の要約 $h_{t-1}$ と結合して新しい要約 $h_t$ を作るのである．したがって RNN は，可変長系列を固定次元の状態遷移として扱う動的システムとみなせる．重み行列が時刻をまたいで共有されるため，系列長が変わっても同じモデルを適用できる点も重要である．

学習は時間方向に展開した計算グラフに対する誤差逆伝播，すなわち BPTT(backpropagation through time)\cite{Werbos1990}によって行う．時刻ごとの損失を $\mathcal L_t$ とすると，総損失 $\mathcal L=\sum_t \mathcal L_t$ に対するある時刻 $t$ の状態勾配は
\begin{equation}
\frac{\partial \mathcal L}{\partial h_t}
= \frac{\partial \mathcal L_t}{\partial h_t}
+ \left(\frac{\partial h_{t+1}}{\partial h_t}\right)^\top
\frac{\partial \mathcal L}{\partial h_{t+1}}
\end{equation}
という再帰式に従う．ここで $\partial h_{t+1}/\partial h_t$ は時刻をまたぐヤコビアンであり，長い系列ではそれらの積が勾配を支配する．このため，固有値の大きさや活性化関数の導関数によって，勾配消失や勾配爆発が生じやすい．

LSTM(long short-term memory)\cite{Hochreiter1997}や GRU(gated recurrent unit)\cite{Cho2014}が必要になったのは，この問題を緩和するためである．たとえば LSTM では，セル状態 $c_t$ を
\begin{equation}
c_t = f_t\odot c_{t-1} + i_t\odot \tilde c_t,
\qquad
h_t = o_t\odot \tanh(c_t)
\end{equation}
と更新する．ここで $f_t,i_t,o_t$ はそれぞれ忘却ゲート，入力ゲート，出力ゲート，$\tilde c_t$ は候補セル状態であり，いずれも時刻 $t$ の入力と前時刻状態から計算される．重要なのは，セル状態が加法的な経路を持つため，勾配が長距離にわたって保たれやすいことである．GRU も同様に，より簡潔なゲーティングによって情報保持と更新を制御する．

RNN は長らく系列モデリングの中心にあったが，固定長の隠れ状態だけで長い文脈を圧縮することには限界がある．Encoder-Decoder 構造は，入力系列を文脈ベクトルに写し，それをもとに出力系列を生成することで機械翻訳などに成功したが，長文では情報ボトルネックが顕著になった．これを緩和したのが Attention\cite{Bahdanau2015} であり，各出力時刻が入力系列全体の関連部分を直接参照できるようにした．この発想がさらに発展し，再帰そのものを不要にしたのが Transformer\cite{Vaswani2017} である．次節では，その構造を詳しく見る．

\subsection{Transformer}
Transformer を学ぶ必要があるのは，現代の大規模言語モデルや多くの生成モデルの中核がこの構造にあるからである．RNN は系列依存を表現できるが，計算が逐次的で並列化しにくく，遠く離れた要素間の依存を長い経路で伝えなければならない．この逐次性と長距離依存性の限界を克服するために，Transformer は系列内の要素同士を直接結び付ける Self-Attention を採用した．これにより，系列長に対する計算構造と表現の仕方が根本的に変わった．

長さ $T$ の系列を，各トークンの埋め込みを並べた $T\times d$ の実数行列 $H$ で表す．ここで $T$ は系列長，$d$ は埋め込み次元である．Self-Attention では，$H$ からクエリ $Q$，キー $K$，バリュー $V$ を
\begin{equation}
Q=HW^Q,\qquad K=HW^K,\qquad V=HW^V
\end{equation}
により作る．$W^Q,W^K,W^V$ は学習される重み行列である．各行はそれぞれ一つの位置に対応し，ある位置が他の位置をどれだけ参照すべきかを，クエリとキーの内積で決める．Scaled Dot-Product Attention は
\begin{equation}
\mathrm{Attention}(Q,K,V)
= \mathrm{softmax}\left(\frac{QK^\top}{\sqrt{d_k}}\right)V
\end{equation}
で定義される．ここで $d_k$ はクエリとキーの次元，softmax は各行ごとに適用されて注意重みを確率分布に変換する．$\sqrt{d_k}$ で割るのは，内積の分散が大きくなりすぎて softmax が飽和するのを防ぐためである．

この式の意味は，各位置が他のすべての位置に対して重み付き平均をとり，新しい表現を得るということである．RNN では時刻 $t$ の情報が時刻 $t+1,t+2,\dots$ を経由して遠方へ伝わるのに対し，Self-Attention では任意の二位置間が一回の演算で直接相互作用できる．したがって，長距離依存性に対する経路長は理論的に短くなり，勾配伝播も容易になる．また，各位置の計算を同時に行えるため，訓練時の並列化効率が高い．

Multi-Head Attention は，この相互作用を単一の重み付けに限定せず，異なる関係を並列に学習するための仕組みである．ヘッド数を $M$ とすると，第 $m$ ヘッドは
\begin{equation}
\mathrm{head}_m
= \mathrm{Attention}(HW_m^Q,HW_m^K,HW_m^V)
\end{equation}
で与えられ，全体の出力は
\begin{equation}
\mathrm{MHA}(H)
= \mathrm{Concat}(\mathrm{head}_1,\dots,\mathrm{head}_M)W^O
\end{equation}
となる．ここで $W_m^Q,W_m^K,W_m^V$ は第 $m$ ヘッド専用の射影行列，$W^O$ は結合後の出力射影，$\mathrm{Concat}$ は特徴次元方向の連結を表す．複数ヘッドを用いることで，あるヘッドは局所的対応，別のヘッドは長距離依存，さらに別のヘッドは構文的関係といった異なるパターンを同時に捉えうる．

Transformer には再帰がないため，位置情報を別途与えなければ系列の順序が失われる．これを補うのが Position Encoding である．代表的な正弦波型位置符号化では，位置 $t$，次元添字 $i$ に対して
\begin{equation}
p_{t,2i} = \sin\left(\frac{t}{10000^{2i/d}}\right),
\qquad
p_{t,2i+1} = \cos\left(\frac{t}{10000^{2i/d}}\right)
\end{equation}
を定め，入力埋め込みに加える．これにより，モデルは相対的・絶対的な位置関係を連続的な表現として受け取ることができる．実際には学習可能な位置埋め込みが使われることも多いが，要点は「順序を表現しない自己注意機構に，系列順序の手がかりを与える」ことである．

RNN と比較した Transformer の利点は，並列化しやすく，長距離依存を扱いやすい点にある．他方で，全位置対の相互作用を計算するため，Self-Attention の計算量とメモリ量は系列長 $T$ に対して概ね $O(T^2)$ で増大する．このため非常に長い系列では効率上の課題が残り，疎注意や線形注意など多くの改良が提案されている．それでも，自然言語処理を中心に Transformer が支配的となったのは，この計算コストを払ってなお得られる表現力と学習効率の利得が大きかったからである．

自己回帰モデル(autoregressive model)との関係も重要である．デコーダ型 Transformer では，未来のトークンを見ないように因果マスク(causal mask)を課し，
\begin{equation}
p_\theta(x_{1:T})
= \prod_{t=1}^T p_\theta(x_t\mid x_{1:t-1})
\end{equation}
という系列分布を学習する．ここで各条件付き分布は，Self-Attention を通じて過去文脈全体に基づいて計算される．この枠組みを巨大データ上で事前学習したものが，大規模言語モデル(large language model; LLM)の基本形である．したがって Transformer は，本章の終点であると同時に，後に学ぶ生成モデルや強化学習と結び付く出発点でもある．

\subsection*{まとめ}
本章では，機械学習を「確率分布の学習と最適化の問題」として捉える立場から，基礎概念を順に整理した．まず歴史的には，パターン認識から統計的学習，ニューラルネットワーク，深層学習へと進む流れを，表現力，計算資源，データ量，最適化という観点で読み直した．次に，教師あり学習の基本設定として，データ集合，モデル，パラメータ，経験リスク，汎化誤差，条件付き分布の対応を定式化した．そこから，損失関数が最尤推定や MAP 推定から自然に導かれること，深層学習の最適化が一般に非凸であることを確認した．

さらに，実際の学習を支える計算原理として，ミニバッチに基づく SGD と，連鎖律に基づく誤差逆伝播法を導入した．これにより，大規模データと多数のパラメータを持つモデルを現実的に学習できる理由が明確になった．そのうえで，ニューラルネットワークの基本構造を，線形変換，非線形活性化，出力層の組合せとして整理し，過学習と汎化の問題に対して正則化が果たす役割を考えた．最後に，画像の局所性を活かす CNN，系列の時間依存を表す RNN，長距離依存と並列化を両立する Transformer を順に見たことで，アーキテクチャとはデータ構造に応じた帰納バイアスの実装であることが理解できたはずである．これらの基礎は，強化学習や生成モデルへ進む際にも，その背後にある確率的定式化と最適化の考え方を見失わないための足場となる．
